\documentclass[11pt]{article}

\usepackage[margin=1in]{geometry}
\usepackage[T1]{fontenc}
\usepackage[utf8]{inputenc}
\usepackage{lmodern}
\usepackage[hidelinks]{hyperref}

\title{Artifact Description}
\date{}

\begin{document}

\maketitle

\section{Abstract}

This artifact description provides the information needed to build the software
environment and launch the core experiments from the SC26 paper
\emph{``Printing in a JIFFY: a Parallel-in-Time Heat Transfer Solver for
Additive Manufacturing.''} The goal of this appendix is to document the
practical steps required to reproduce the results reported in the paper's
empirical evaluation.

We begin by describing the target machine configuration, required software
dependencies, loaded modules, and programming environment. We then organize the
artifact into three parts. Part~1 focuses on generating the graph-partitioning
schedule used by the parallel execution framework. Part~2 covers the
GPU-accelerated linear solver for heat transfer on a moving computational
domain, inherited from the HERMES project. Part~3 presents the multi-GPU
parallel solver, which combines the scheduling machinery from Part~1 with the
numerical infrastructure from Part~2 to reproduce the end-to-end parallel
experiments.

\section{Artifact Meta Information}

The artifact is implemented primarily in Python and targets NVIDIA GPU-based
systems. The core numerical kernels and linear algebra routines rely on CUDA,
CuPy, Numba-CUDA, SciPy, and NumPy. Some path-generation modes additionally
require \texttt{scikit-image}.

The experiments reported in this artifact are designed for the following
software and hardware environment:
\begin{itemize}
\item Platform: TACC Vista, a GPU-enabled high-performance computing system at
the Texas Advanced Computing Center.
\item Compute node: NVIDIA Grace Hopper GH200 nodes, which serve as the target
platform for both the single-GPU solver and the multi-GPU parallel runs.
\item Programming language: Python~3.11.0.
\item Required system modules: \texttt{gcc/15.1.0}, \texttt{openmpi/5.0.5},
and \texttt{cuda/12.5}.
\item Core GPU and scientific packages: CUDA~12.5, CuPy
(\texttt{cupy-cuda12x}~13.6.0), Numba-CUDA (\texttt{numba}~0.63.1),
SciPy~1.17.0, and NumPy~2.3.5.
\item Utility packages: Matplotlib, PyYAML, and \texttt{psutil}.
\item Optional package: \texttt{scikit-image}, required for image-based path
configurations.
\item Parallel runtime: OpenMPI~5.0.5 for multi-GPU experiments.
\end{itemize}

\section{Partition Schedule Generation}

Part~1 of the artifact corresponds to the planning-only workflow implemented in
\texttt{src/hermes/scripts/segment\_correction/plan\_only.py}. This workflow
constructs a graph-partition schedule without launching the full parallel
simulation. Its role in the artifact is to generate the path segments and the
runtime partition used later by the parallel solver.

The inputs to this stage are:
\begin{itemize}
\item A simulation configuration file, such as \texttt{configs/sim\_ex1.ini},
which specifies the physical and solver parameters.
\item A path configuration file, such as \texttt{configs/fast\_heat.ini}, which
defines the laser trajectory.
\item Planning parameters including \texttt{world\_size},
\texttt{planner\_mode}, \texttt{dt\_us}, and \texttt{correction\_weight}.
\end{itemize}

The main workflow is straightforward. First, the planner reads the path
configuration and generates a laser trajectory. For image-based inputs, this
means reading a binary image and converting its boundary into a scan path with
an interior zigzag fill pattern. This stage is handled by
\texttt{build\_path\_sections\_nd\_from\_ini()} in
\texttt{src/hermes/laser\_path/path\_loader.py} and, for picture-mode inputs,
by \texttt{build\_picture\_nd()} in
\texttt{src/hermes/laser\_path/path\_builders.py}.

Second, the generated path is discretized into laser segments. This is done by
sampling the trajectory and converting it into a time-discretized laser scan
sequence, followed by \texttt{build\_segments\_from\_program()} in
\texttt{src/hermes/pipelines/ss\_builder.py}. In the current artifact, the
planning granularity is the segment itself.

Third, the planner builds a dependency DAG over the resulting ordered segment
sequence. This is performed by \texttt{compute\_dag\_and\_components()} in
\texttt{src/hermes/pipelines/components.py}, which constructs dependency edges
and identifies the structure needed for later partitioning.

Finally, the DAG is partitioned across the target number of ranks. The main
entry point is \texttt{build\_runtime\_plan()} in
\texttt{src/hermes/scheduling/planning.py}, which invokes the partitioning
logic in \texttt{src/hermes/scheduling/\_partitioned.py}. Under the default
\texttt{exact\_dp} mode, the planner assigns contiguous intervals of the ordered
segment sequence to ranks while minimizing the predicted maximum workload. The
result is written to \texttt{planning\_summary.json}, which is then used by the
parallel execution stage.

\section{HERMES Outer-Level Solver with Moving Computational Domain}

Part~2 of the artifact is inherited from the open-source HERMES project
(\url{https://github.com/aydinalperen7/hermes-gpu-heat}). In this artifact,
we reuse the HERMES outer-level domain and outer-level GPU heat solver with a
moving computational domain. Its role in the artifact is to provide the
single-GPU numerical solver that advances the temperature field while the laser
moves along the prescribed scan path.

The main solver implementation is located in
\texttt{src/hermes/scripts/outer\_solver.py}. This solver operates on the
Level-3 outer domain and updates the thermal field as the computational domain
translates with the laser. This moving-domain design allows the solver to focus
computational effort on the active region around the heat source instead of
evolving a large fixed domain over the entire build area.

The solver supports two execution modes: \texttt{legacy} and \texttt{fused}.
The \texttt{legacy} mode uses a more conventional GPU implementation of the
linear solve, while the \texttt{fused} mode uses custom CUDA kernels to combine
multiple operations and reduce overhead. In the current artifact, the default
mode is \texttt{fused}, and it is the mode used in the main experiments.

\section{Multi-GPU Parallel Solver}

Part~3 combines the partition schedule from Part~1 with the outer-level solver
from Part~2 to form the full multi-GPU parallel execution workflow. Its role in
the artifact is to combine the partition and schedule with the underlying
single-GPU solver to produce the parallel-in-time execution. The main entry
point is \texttt{src/hermes/scripts/segment\_correction/main.py}, with the core
parallel execution logic in
\texttt{src/hermes/scripts/segment\_correction/core.py}. The parallel execution
model is based on MPI.

At the beginning of the run, rank~0 generates the partition and schedule, then
broadcasts the runtime path definitions and rank assignments to all other
processes through MPI. After that, each rank executes the components assigned
to it.

For each assigned component, a rank first runs the base simulation with the
laser source turned on over its local path segment. It then computes the
corresponding source-off correction needed for downstream components. When that
correction is required by a component owned by another rank, the correction data
is communicated through MPI to the target rank. The receiving rank then applies
superposition by combining the received correction with its local base state.

In this way, MPI serves as the communication layer for distributing the
schedule, exchanging inter-rank correction data, and coordinating the overall
multi-GPU execution. Conceptually, this stage augments the HERMES solver with
rank-level scheduling and correction exchange to produce the final
parallel-in-time solver.

For reproducibility, the bash scripts \texttt{run\_mpi\_sbatch.sh} and
\texttt{run\_mps\_sbatch.sh} in the \texttt{Jiffy} root directory can be used
to launch the multi-GPU experiments and reproduce the strong-scaling results.

\end{document}
